% SPDX-FileCopyrightText: Copyright (c) 2026 NVIDIA CORPORATION & AFFILIATES. All rights reserved.
% SPDX-License-Identifier: Apache-2.0
%
% ROLE: place the work among (1) KV-cache-aware routing heuristics, (2) learned routers for LLM
%   serving, (3) learned systems policies trained in simulation and black-box policy search,
%   (4) discrete-choice models with context effects, (5) goodput, SLO and evaluation methodology.
%   Each paragraph says what we take from the area and where we differ; it is not a survey.
% EVIDENCE: CR/literature/LESSONS.md (adversarially verified lessons; "VERIFIED (pdf p.N)" lines),
%   CR/literature/notes/*.md (scout notes with section and page pointers), CR/literature/BIBLIOGRAPHY.md
%   (read depth). Claims not in LESSONS.md were re-checked against CR/literature/txt/ or
%   pdftotext of CR/literature/pdfs/ on 2026-10-03 (marked "re-checked" below). Sources read only as
%   an abstract or web page are cited for what that abstract or page says, nothing more.
% STATUS: final for this draft. Restructured 2026-10-05: Lodestar's 80% threshold moved here from the
%   introduction; the set-dependent-models paragraph compressed with every citation kept. Round-1 fix
%   (E1): "every result outside the live check is simulated".

\section{Related work}
\label{sec:related}

Five lines of work meet in this study. For each we say what we take from it and where we differ.

\paragraph{KV-cache-aware routing.}
Most of our baselines come from this line (\cref{sec:baselines}), and the published rules weigh
prefix reuse against load in different ways.
% src: CR/literature/notes/llm-routing.md 2.4 (Preble sec. 3.2, p.5); re-checked txt/llm-routing/srivatsa2024-preble.txt l.67-80
Preble sends a request to the GPU holding its prefix when the shared part outweighs the rest of
the prompt, and otherwise to the GPU with the lowest cost, which counts recent load, the eviction
the request would cause and the prefill of its missed tokens \citep{srivatsa2025preble}.
% src: LESSONS.md LR-07 (LMetric sec. 4.4: 0.7 ChatBot, 0.55 API; pdf p.6, VERIFIED); LR-04 (P-token, pdf p.2, p.8-9)
LMetric replaces weighted sums with the product of a worker's new prefill tokens and its batch
size, arguing that the best weight of a linear score changes with the workload: 0.7 on one of
its traces and 0.55 on another \citep{zhang2026lmetric}.
% src: LESSONS.md LR-07 (DualMap sec. 2, pdf p.5-6, VERIFIED)
DualMap restricts each prefix to two hashed candidates; it argues that picking the best of all
workers amounts to a power-of-$\Nw$-choices rule that scatters shared prefixes, and that a
per-request minimum-TTFT rule oscillates between cache-aware and load-aware decisions
\citep{yuan2026dualmap}.
% src: LESSONS.md LR-07 (Lodestar sec. 4.1: k = 2 consistent-hash filter when cluster memory utilization
%      exceeds 80%, pdf p.6, VERIFIED)
Lodestar applies a similar two-choice consistent-hash filter, but only while cluster memory
utilization is above a threshold, 80\% in its evaluation \citep{lim2026lodestar}.
% src: LESSONS.md LR-12 (CacheRoute Table 2, pdf p.4, VERIFIED); re-checked txt/llm-routing/cheng2026-cacheroute.txt
%      l.579-586 ("A cache-hit increase is not enough; the plan is enabled only when p99 or capacity improves")
CacheRoute periodically plans which replicas serve the most popular prefixes, and gates each plan
on a shadow replay because a higher cache-hit rate alone did not guarantee lower tail latency or
higher capacity \citep{cheng2026cacheroute}.
% src: BIBLIOGRAPHY.md (llm-d blog, read on the web); the port's provenance in llm_d/precise_prefix.rs
llm-d scores workers with a prefix scorer fed by an index of the engines' KV events
\citep{llmd2025kvcache}.
For agent traffic,
% src: LESSONS.md LR-07 (SMetric Fig. 13, pdf p.8, VERIFIED)
SMetric sends a session's follow-up turns to the worker with the most cached context unless that
worker's estimated load would break a slack multiple of the TTFT SLO \citep{wang2026smetric};
% src: BIBLIOGRAPHY.md (Autellix, read on the web, sec. 4.3); CR/literature/notes/llm-routing.md 2.7
Autellix keeps an agent program's long calls on the program's engine and sends short calls to
the least-used engine \citep{luo2025autellix};
% src: LESSONS.md LR-14 (ThunderAgent sec. 3.2, pdf p.4-5, VERIFIED)
and ThunderAgent schedules whole programs, observing that a router-side prefix tree that
approximates worker caches goes stale under eviction thrashing \citep{kang2026thunderagent}.
% src: CR/literature/notes/llm-routing.md 2.7 (Nixon et al. sec. 7.2, read on the web)
A workload study simulated cache-first, load-first and sticky routing at several instance counts
and tracked hit ratio, per-session replication and load imbalance \citep{nixon2026year}.

% src: LESSONS.md LR-07 (SMetric sec. 4.1, pdf p.6: 920 vs 1,437 TPS within SLO, reuse 45% vs 74%,
%      no global tier; VERIFIED); LR-04 (advantage "except for setups without any global KV$ store", pdf p.10)
These rules disagree about which signal should dominate. On an agentic workload without a global
KV tier, which is the setting of our replay, SMetric measured LMetric serving 36\% less
SLO-meeting throughput than a tuned linear cache-plus-load score, with prefix reuse of 45\%
against 74\% \citep{wang2026smetric}. We do not propose another heuristic. We tune these rules on
equal terms, and the learned policy's features draw on their signals: uncached prefill, active
prefill, batch size, KV utilization and session affinity, and, in feature set \code{v2}, LMetric's
log-product and a hashed home-worker indicator in the spirit of DualMap
(\cref{tab:features,tab:features-v2}).
% src: FEATURES.md "Feature set v2" (log_ptok, log_bs: LR-07 LMetric; hash_home: LR-07 Lodestar, DualMap)

\paragraph{Learned routers for LLM serving.}
Learned routers differ in what they learn and in where they learn it.
% src: CR/literature/notes/llm-routing.md 1.3 (Jain et al. sec. 5.3); re-checked
%      txt/llm-routing/jain2024-intelligent-router.txt l.895-898 ("a in {0, ..., m}")
Jain et al.\ train a reinforcement-learning router with a heuristic-guided reward; its action
picks one of $m$ instances or none, so the network's shape is tied to the instance count
\citep{jain2024intelligent}.
% src: LESSONS.md LR-06 (Lodestar sec. 4.1, pdf p.5, VERIFIED: shared parameters, no instance identity);
%      re-checked txt/llm-routing/lim2026-lodestar.txt l.488-505 (circular dependency, Fig. 3a)
Lodestar learns online an MLP that predicts each instance's reward from per-instance features,
with parameters shared across instances and no instance identity, so it does not depend on the
instance count. It also shows that a latency predictor trained offline on logs of other routers
loses its ranking signal once it routes, because the routing policy changes the states the
predictor is evaluated on \citep{lim2026lodestar}.
% src: LESSONS.md LR-13 (GORGO sec. 5, App. D, pdf p.7-8, VERIFIED)
GORGO tunes the weights of an additive routing cost with an evolution strategy; optimizing p95
TTFT alone drove the queue weight to zero and sent every request to the closest replica
\citep{toniolo2026gorgo}.
% src: BIBLIOGRAPHY.md (Wu et al., ICLR 2026, key sections); title-level claim only
Wu et al.\ analyze KV eviction and learning-based load balancing together
\citep{wu2026randomization}.
% src: BIBLIOGRAPHY.md (Jha et al. 2024, RouteBalance 2026: abstract only, "tangential")
Routers that choose among models or service-quality levels, rather than among replicas of one
model, address a different decision \citep{jha2024learned,da2026routebalance}.
% src: LESSONS.md LR-02 (CtR: hand-set weights, only constants fitted, pdf p.4, p.6), LR-04 (sec. V-B:
%      0.819 vs 0.864, pdf p.3-5), LR-14 (sec. VI: residual 0.068, "does not reproduce the measured
%      winner on any workload", pdf p.5-6), LR-13/LR-12 (sec. V-E: 0.292 vs RR 0.680, pdf p.4-5); all VERIFIED
The closest study in method is Calibrate-then-Route \citep{tumkur2026calibrate}, a router with
hand-set weights and fitted constants that was developed in a simulator and then measured on GPUs
in disaggregated serving. With simulator-derived constants it lost 4.5 goodput points to the same
router with measured constants (0.819 against 0.864). After re-fitting, the simulator matched
measured goodput to a mean absolute residual of 0.068 but did not reproduce the measured winner
on any workload. In a starved pool, greedy cost minimization collapsed to 0.292 goodput while
round-robin held 0.680.
% src: LESSONS.md LR-09 (AgentServeSim, causal release, pdf p.1-2, p.4) and LR-02 (sec. 5 Table 2:
%      0.5% KV retention, 2.8% scheduling, pdf p.8); both VERIFIED
AgentServeSim releases agent turns causally in simulation and searches policies with an
LLM-driven evolutionary loop; its best evolved policies improved mean job completion time by
0.5\% for KV retention and 2.8\% for scheduling over hand-written starting policies
\citep{rajib2026agentservesim}.
Our policy differs from these in three ways. It is trained on-policy, since every candidate is
scored by replaying it, which avoids the logged-policy mismatch Lodestar describes, and it is
frozen afterwards. Its objective counts a request as good only if both its inter-token latency
and its end-to-end slowdown are within bounds, where GORGO's objective was p95 TTFT alone. And
Calibrate-then-Route's findings are why ranking
stability under router-state lag and on GPUs is part of the protocol
(\cref{sec:robustness,sec:live}).
% src: CR/literature/notes/llm-routing.md 4.3 ("None of the learned routers report training on some N and
%      testing on unseen N, except Lodestar's architectural N-independence")
Among the learned routers we read, none trains at some worker counts and tests at unseen ones;
Lodestar's shared predictor is independent of the instance count by construction.

\paragraph{Learned systems policies and black-box search.}
% src: LESSONS.md LR-03 (Decima sec. 5.3 fixed arrival sequences, pdf p.7, p.10; Mao ICLR'19, pdf p.3, p.5),
%      LR-10 (Decima App. I Table 3: 630 vs 610 s, pdf p.19), LR-15 (Decima Table 2: 104.8 s vs 91.2 s,
%      pdf p.11); all VERIFIED
Decima learns cluster scheduling with reinforcement learning in a simulator, and it fixes the
job-arrival sequence across episodes to cut the variance that stochastic arrivals add
\citep{mao2019decima,mao2019inputdriven}. It shows both sides of transfer: an agent trained on a
cluster ten times smaller was 3\% worse in average job completion time (630 against 610 s),
while one trained on a mismatched load did worse than the tuned heuristic (104.8 against 91.2 s)
\citep{mao2019decima}.
% src: re-checked txt/learned-systems-policies/mao2017-pensieve.txt l.298-360 (slow-start-restart disabled)
Pensieve trains a bitrate controller in a chunk-level simulator and disabled TCP slow-start
restart on its video server so that deployment matched a simulator assumption
\citep{mao2017pensieve}.
% src: LESSONS.md LR-07 (Park sec. 3.1 bootstrapping, pdf p.4) and LR-13 (Park sec. 3.3 reservation,
%      pdf p.5); VERIFIED
Park catalogs the difficulties of learning systems policies; one of its load balancers learned to
keep a server free for small jobs and was not robust when the job mix changed, and it suggests
bootstrapping search from existing policies \citep{mao2019park}, as M1's heuristic-derived
starting points do (\cref{sec:restarts}).
% src: LESSONS.md LR-14 (Puffer sec. 5.2, pdf p.9-10), LR-09 (CausalSim, pdf p.1-3), LR-11 (Genet,
%      pdf p.3-4), LR-14 (Peng Tables II-III, pdf p.7), LR-13 (Lehman, pdf p.9-10); all VERIFIED;
%      Puffer abstract re-checked txt/learned-systems-policies/yan2020-puffer.txt l.22-32
The cautions are empirical. Puffer's randomized trial found it difficult for sophisticated or
learned bitrate controllers to outperform a simple buffer-based scheme despite good results in
emulation, and a controller whose predictor was trained in emulation performed poorly in
deployment
\citep{yan2020puffer}. Trace-driven simulation is biased when the policy changes the traces it
replays \citep{alomar2023causalsim}. Learned policies that win on average can still lose to the
baselines in a substantial fraction of environments \citep{xia2022genet}. Randomizing latency
and observation noise was necessary for simulation-trained robot control to transfer
\citep{peng2018sim2real}. And search exploits simulator bugs \citep{lehman2020creativity}. These
findings motivate the closed-loop and agent-lane cells, the per-cell winner map, the timing and
lag perturbations, the concentration audit and the live check (\cref{sec:eval}).

% src: re-checked txt/learned-systems-policies/mania2018-ars.txt l.17, l.84 (static linear policies);
%      LESSONS.md LR-05 (ARS sec. 3.2 state normalization, pdf p.7-8; Hansen tutorial sec. 5, pdf p.24),
%      LR-03 (PEGASUS, pdf p.3-7; UH-CMA-ES, pdf p.4, p.8-9); re-checked
%      txt/learned-systems-policies/salimans2017-es.txt l.282-290, l.450 (variance independent of T)
On method, static linear policies trained by random search matched reinforcement learning's
sample efficiency on the MuJoCo benchmarks, given normalized inputs \citep{mania2018ars};
evolution strategies' gradient estimates do not depend on episode length, which suits long
horizons with delayed reward \citep{salimans2017es}. We use CMA-ES \citep{hansen2016tutorial},
score every candidate of a generation on the same random scenarios, as PEGASUS fixes scenarios to
make policy evaluation deterministic \citep{ng2000pegasus}, and borrow UH-CMA-ES's rank-change
measurement as a diagnostic, without its treatment of the measured uncertainty
\citep{hansen2009uh}.
% src: LESSONS.md LR-14 (Mitzenmacher, pdf p.2; Tahir sec. 1, pdf p.1-2), LR-12 (van der Boor sec. 1, 2.4,
%      pdf p.4, p.7, p.11-12); VERIFIED
Load-balancing theory frames two risks of our setting: sending work to the least-loaded server
on old load information can significantly hurt performance \citep{mitzenmacher1998old}, and
delayed information makes join-the-shortest-queue dispatch herd \citep{tahir2022meanfield}; how balancing rules compare
also changes with the number of servers and the load regime \citep{vanderboor2022survey}.

\paragraph{Discrete choice with context effects.}
With rank 0, \lc{} is a conditional logit, the base case of the generalized extreme-value family
\citep{train2009gev}: one coefficient vector scores every candidate, so the same parameters serve any worker count, and at $\temp = 0$
the policy takes the logit's most likely choice.
% src: re-checked pdftotext CR/literature/pdfs/dcm-context/maragheh2018-halo-mnl.pdf (sec. 1: two-way
%      interaction terms between items); BIBLIOGRAPHY.md ("design confirmation (no worker-ID terms)")
Halo models add pairwise terms between specific items \citep{maragheh2018halo}; \lc{} has no
worker-identity terms, so its parameters apply unchanged as workers come and go.
% src: re-checked pdftotext CR/literature/pdfs/dcm-context/rosenfeld2020-set-dependent-aggregation.pdf
%      (sec. 1: argmax predictions "are clearly independent of s"; cannot be fixed by more complex F)
A plain logit satisfies the independence of irrelevant alternatives: the preference between two
candidates never depends on which others are present, however rich the item score
\citep{rosenfeld2020setdependent}. M2's context terms and \code{v2}'s set-relative features add
that dependence.
% src: LESSONS.md LR-05 (Tomlinson sec. 3.1-3.2, pdf p.4-5: column-restricted LCL, not DLCL), LR-06
%      (Lemma 4.3, Prop. 4.5, pdf p.4-6), LR-15 (footnote 3, Fig. 2, pdf p.10-11, p.13); VERIFIED
In the linear context logit (LCL) of \citet{tomlinson2021lcl} the coefficient vector shifts
linearly with the mean features of the choice set, $\theta + A\fbar$. M2's pooled form is an LCL
with a low-rank $A$; its named-source form is an LCL over a vector of named set statistics, with
one free column of $A$ per statistic. Neither is their DLCL, a mixture of logits with separate
intercepts. Three of their results shaped our design. Choice depends only on deviations from the set mean, so terms that are constant across
the set cancel. $A$ is identifiable only from at least $\dfeat + 1$ choice sets with affinely
independent mean features. And LCL did not beat the plain logit within a fixed training budget
although it nests it, one reason the ladder adds context in a small step that starts from the
selected M1 and is compared with M1 trained for the same added evaluations (\cref{sec:ladder}).

% src: LESSONS.md LR-05 (Seshadri pdf p.4, p.8; Rosenfeld Fig. 3, pdf p.7-8; Ko & Li eq. 3, pdf p.5),
%      LR-06 (Pfannschmidt pdf p.7, p.11, p.22-23; Webb pdf p.26-30), LR-15 (Aouad pdf p.8; DeepHalo
%      pdf p.5-6; Train pdf p.2-9); VERIFIED
Other set-dependent models sum pairwise context vectors over the set \citep{seshadri2019cdm} or
use set-dependent aggregation \citep{rosenfeld2020setdependent}, pooled set embeddings
\citep{pfannschmidt2022context}, mixtures of random-utility networks \citep{aouad2023rumnet},
self-attention \citep{ko2024attention} or stacked interaction layers \citep{zhang2026deephalo};
divisive normalization, which fits human choices across set sizes from 2 to 12
\citep{webb2021normalization}, inspired \code{v2}'s ratio features. Low-rank choice parameters are
unique only up to a multiplicative factor unless a penalty fixes the scale \citep{ko2024attention};
M2's pooled form shares that freedom (scaling $p_{\cidx}$ up and $q_{\cidx}$ down by one factor
changes nothing), while its named-source form, linear in $p_{\cidx}$, does not. A nested logit
groups alternatives whose unobserved utilities are correlated \citep{train2009gev}, but under a
deterministic argmax its nest parameters have no effect (our inference), so node structure would
enter as nest-level features rather than as a GEV model (\cref{sec:nested}).
% src: re-checked pdftotext CR/literature/pdfs/dcm-context/jain2020-generalization-new-actions.pdf
%      (softmax over available actions; subsampled action sets; validation on held-out action sets,
%      Alg. 1 l.17); LESSONS.md LR-10 (pdf p.4-5, p.8, VERIFIED); LR-06 (Deep Sets, corrected scope)
In reinforcement learning, \citet{jain2020newactions} use a shared utility followed by a softmax
over the available actions, a conditional logit over a variable action set; generalizing it to
unseen action sets took subsampled action sets, entropy regularization and model selection on
held-out action sets. Our analogue of the last is validation at $\Nw = 6$, a worker count
training never sees. The named context sources are request-level quantities or means and maxima
over candidates, a permutation-invariant pooling \citep{zaheer2017deepsets}, chosen so that
duplicating every worker leaves them unchanged.
% src: FEATURES.md "Context sources" (unit test duplicating_the_worker_set_keeps_features_sources_and_choice)

\paragraph{Goodput, SLOs and evaluation methodology.}
% src: re-checked txt/eval-methodology/zhong2024-distserve.txt l.76-82 (goodput = max request rate at the
%      SLO attainment goal); LESSONS.md LR-08 (SLO-scale sweep, pdf p.10, VERIFIED)
DistServe defines goodput as the highest request rate that meets an SLO-attainment target, and
sweeps a single scale factor over its SLOs \citep{zhong2024distserve}. We measure good requests
per second at a fixed load inside a window, and judge a request by its mean inter-token latency
and its end-to-end slowdown against an uncontended latency (\cref{sec:objective}); as our
analogue of DistServe's sweep, the finalists' test runs are re-scored at several SLO scales.
% src: LESSONS.md LR-08 and LR-13 (Wang et al. sec. 4.2 eq. 8 and sec. 1, pdf p.1, p.6, VERIFIED)
Binary SLO goodput rewards giving up on requests that will miss anyway, and a TPOT-style bound
averages out stalls in the middle of a request \citep{wang2024slo}; both are why the audits look
for sacrificed requests and concentrated load (\cref{sec:audits}).
% src: LESSONS.md LR-09 (Schroeder principles (v), (vii), pdf p.10-11, VERIFIED; scope: an analogy for routing)
Open and closed load generators behave differently: scheduling helps a closed system only at
moderate load, and, as a rule of thumb, sessions of more than ten requests make a partly open
system behave as closed \citep{schroeder2006open}; we therefore stratify by load mode.
% src: LESSONS.md LR-14 (Vidur sec. 7.2, pdf p.9, VERIFIED); BIBLIOGRAPHY.md (Abdelfattah et al.:
%      abstract only; "22.2%" and "14 industrial cases" from the abstract, CR/literature/notes/eval-methodology.md sec. 4)
Simulator error grows near capacity: in Vidur, small prediction errors near the capacity point
blow up into large latency errors \citep{agrawal2024vidur}. A load-testing study reports that
handling warm-up properly improved the accuracy of capacity estimates by 22.2\% across 14
industrial cases \citep{abdelfattah2026loadtesting}.
% src: LESSONS.md LR-10, LR-11 (Demsar pdf p.5, p.7; Agarwal pdf p.5-8; Bouthillier pdf p.8-10;
%      Henderson pdf p.5; Dodge pdf p.1-3; Cawley pdf p.1, p.16, p.25); VERIFIED
Our statistics follow the guidance for comparisons over few independent tasks: a Wilcoxon
signed-rank test over independent units \citep{demsar2006}, bootstrap intervals and
interquartile means over tasks \citep{agarwal2021precipice}, a probability-of-improvement effect
size with thresholds for significant and meaningful \citep{bouthillier2021variance}, several
seeds and restarts \citep{henderson2018matters}, budget-aware reporting \citep{dodge2019show}, and
selection on data the final test never sees \citep{cawley2010overfitting}.

\paragraph{Position of this work.}
The study combines four choices that we did not find together in the work above. The policy is a
deployable choice rule, linear in inputs the live router already has, rather than a latency
predictor. Every baseline is tuned with the same budget and selection rule as the learned model,
and the headline comparison is against the validation-selected best of them. Worker counts are
held out: training uses $\Nw \in \{4, 8\}$, and the test adds 2, 6, 16 and 32. And the headline
test is pre-registered over independent trace segments. Two of our findings echo this literature.
The learned policy's advantage over the ported heuristics is largely an LMetric-style
interaction between uncached prefill and load (\cref{sec:res-coefficients,sec:res-fairness}), which
supports LMetric's point that the two should interact, though here as a correction added to the
default's additive cost rather than in place of it; and the first unconstrained
search found the kind of objective gaming GORGO and Lehman et al.\ document, which a sign constraint
on load coefficients removed (\cref{sec:gaming}). The price is that every result outside the live
check is simulated; as Calibrate-then-Route and Puffer show, a simulated ranking need not survive
deployment.
Our live check on six test cells found the simulated ranking intact on GPUs, but not the simulated
gap sizes (\cref{sec:live}).
